% Inactive independent reconstruction, checkpoint601. No claim for all continuous data.
\section{Fredholm theory for a continuous Hermitian form}
\label{sec:continuous-hermitian-fredholm}

A continuous positive Hermitian form still defines a uniformly elliptic
trace operator. Its range is most conveniently described in Sobolev spaces:
continuity of the coefficients suffices for the local $W^{2,p}$ estimate,
while the maximum principle identifies the kernel.

\begin{lemma}\label{lem:continuous-hermitian-fredholm}
Let $\sigma$ be a continuous positive real $(1,1)$-form on a connected
compact complex $n$-fold, and let $\mathrm dV$ be a smooth positive volume.
For $p>2n$, the real operator
\[
 L_\sigma v=\operatorname{tr}_\sigma\ddc v,
 \qquad L_\sigma:W^{2,p}(X)\longrightarrow L^p(\mathrm dV),
\]
is Fredholm of index zero. Its kernel consists of the constants, and
its range has a one-dimensional annihilator in $L^{p'}(\mathrm dV)$,
where $p'=p/(p-1)$.
\end{lemma}
\begin{proof}
Fix a smooth positive Hermitian form $\eta$ and put
$L_s=(1-s)L_\eta+sL_\sigma$, $0\leq s\leq1$. On a finite holomorphic
cover these operators have uniformly positive real principal matrices
and a common modulus of continuity. They have no lower-order terms in
these coordinates. Write $d=2n$ for the real dimension.

\emph{Local estimates.} For a constant positive real matrix $a_0$, the
symbols $\xi_i\xi_j/(1+a_0^{kl}\xi_k\xi_l)$ satisfy the order-zero
multiplier estimates. Consequently
\[
 \|D^2z\|_p\leq C\bigl(\|a_0^{ij}D_{ij}z\|_p+\|z\|_p\bigr)
 \qquad(z\in C_c^\infty).
\]
Freeze the principal matrix of $L_s$ at a point and shrink the ball until
the coefficient oscillation times this constant is less than $1/2$.
The oscillation term can then be absorbed on the left. Differentiating
a cutoff adds first derivatives; the interpolation inequality
$\|Dz\|_p\leq\varepsilon\|D^2z\|_p+C_\varepsilon\|z\|_p$
absorbs them. This inequality follows by comparing $Dz$ with its
mollification at scale $\varepsilon$ and estimating the derivative of
the mollified function by $C\varepsilon^{-1}\|z\|_p$.
A partition of unity and smooth density give, uniformly in $s$,
\begin{equation}\label{eq:continuous-hermitian-wtwo}
 \|v\|_{W^{2,p}}\leq C\bigl(\|L_sv\|_p+\|v\|_p\bigr).
\end{equation}

\emph{Maximum principles.} We give the local estimate needed to pass
comparison to $W^{2,p}$. Let $L=a^{ij}D_{ij}$ on a bounded smooth
coordinate domain $D$, with $a\geq\vartheta I$, and let $z$ be smooth
with nonnegative boundary values and minimum $-M<0$. Set
$R=\operatorname{diam}D$. Every vector $\xi$ with $|\xi|<M/(2R)$ is
the gradient at an interior lower supporting point of $z$: minimize
$z(x)-\xi\cdot(x-x_0)$, where $z(x_0)=-M$. Its boundary values exceed
$-M/2$, whereas its value at $x_0$ is $-M$. At a supporting point
$D^2z\geq0$, and the arithmetic--geometric mean inequality gives
\[
 \det D^2z\leq
 \left(\frac{\operatorname{tr}(aD^2z)}{d\vartheta}\right)^d.
\]
The area formula for the gradient, restricted to these supporting
points, therefore bounds the volume of the slope ball by
$(d\vartheta)^{-d}\int_D((Lz)^+)^d$. Thus
\begin{equation}\label{eq:continuous-hermitian-local-minimum}
 -\min_Dz\leq C_d\vartheta^{-1}R\|(Lz)^+\|_{L^d(D)}.
\end{equation}
This calculation uses no derivative of $a$. Approximate a $W^{2,p}$
function smoothly on the closure of a smaller smooth domain. The
approximation converges in $C^1$, and its $L$ errors converge in $L^d$
because $p>d$. Correct the vanishing boundary error in
\eqref{eq:continuous-hermitian-local-minimum}. Passage to the limit
shows that $Lv\leq0$ and $v|_{\partial D}\geq0$ imply $v\geq0$.

The same comparison proves the strong maximum principle. If
$v\in W^{2,p}$, $Lv\geq0$, and $v$ has an interior maximum $M$, put
$w=M-v\geq0$. If $w$ is positive somewhere and zero somewhere, a small
ball contained in $\{w>0\}$ and tangent to its zero set can be chosen
inside a coordinate chart. On the outer half-annulus of this ball,
\[
 \zeta(x)=\exp(-k|x-c|^2)-\exp(-kr^2)
\]
is positive and has $L\zeta>0$ for sufficiently large $k$, by uniform
ellipticity. A small positive multiple of $\zeta$ is bounded above by
$w$ on both boundary spheres, hence throughout the annulus by comparison.
At the tangent point this forces a positive inward derivative of $w$,
contradicting $Dw=0$. Connectedness proves that $v$ is constant.
Applied to $L_s$ on $X$, it gives $\ker L_s=\mathbb R$. If
$(L_s-1)v=0$ and $v$ has a positive maximum, apply the strong maximum
principle on its positive component, where $L_sv=v>0$. Constancy
contradicts that equation. Applying the argument to $-v$ rules out a
negative minimum. Hence $\ker(L_s-1)=0$.

\emph{Index and range.} Compactness of $W^{2,p}\hookrightarrow L^p$
and \eqref{eq:continuous-hermitian-wtwo} now give
\begin{equation}\label{eq:continuous-hermitian-inverse}
 \|v\|_{W^{2,p}}\leq C\|(L_s-1)v\|_p
\end{equation}
uniformly in $s$. Otherwise normalize a contradicting sequence in
$W^{2,p}$, pass to $s_j\to s$ and a strongly $L^p$ convergent subsequence,
and apply \eqref{eq:continuous-hermitian-wtwo} to differences. Operator
norm continuity in $s$ gives strong $W^{2,p}$ convergence to a nonzero
member of $\ker(L_s-1)$, which is impossible.
The smooth operator $L_\eta-1$ has index zero and trivial kernel, hence
is an isomorphism, by smooth elliptic theory. The method of continuity
and \eqref{eq:continuous-hermitian-inverse} show that every $L_s-1$ is
an isomorphism. Finally, $L_\sigma$ differs from $L_\sigma-1$ by the
compact inclusion $W^{2,p}\hookrightarrow L^p$, so it is Fredholm of
index zero. Its one-dimensional kernel makes its range have codimension
one. $L^p$ duality gives the annihilator assertion.
\end{proof}

\subsection{Continuous positive complex Hessians}
\label{subsec:continuous-hessian-adjoints}

The smooth recovery theorem extends to a bounded solution whose complex
Hessian is continuous and everywhere positive. In this case one can
vary the potential in both directions, and the resulting densities
force every supporting probability to annihilate the same elliptic
operator. The Sobolev range theorem then gives uniqueness.

\begin{theorem}\label{thm:continuous-hessian-adjoint-uniqueness}
Let $n\geq2$, let $\theta=\omega$ be a Hermitian form, and let $\varphi=0$.
Suppose $F\in C(X)$ and a bounded solution $u\in\PSH(X,\omega)$ of
\eqref{eq:relative-normalization} has distributional complex Hessian
represented by a continuous form
\[
 \sigma=\omega+\ddc u>0
\]
everywhere on $X$. Then $\mathcal A_0(F)$ consists of one probability
$\nu_F$. It satisfies
\[
 \int_X L_\sigma h\,\mathrm d\nu_F=0\quad(h\in C^\infty(X)),
 \qquad
 DB_0(F)[H]=-\int_XH\,\mathrm d\nu_F\quad(H\in C(X)).
\]
The derivative is taken in the G\^ateaux sense.
\end{theorem}
\begin{proof}
The continuous trace of $\ddc u$ is the Euclidean Laplacian of $u$, up
to the fixed normalization. Local constant-coefficient elliptic
regularity gives $u\in W^{2,p}$ for every finite $p$. Indeed, subtract
the Newtonian potential of a cutoff of this trace. Its second
high-frequency derivatives are order-zero multipliers, and the
low-frequency kernels are locally bounded; the remaining distribution
is harmonic. Sobolev embedding gives a $C^{1,\alpha}$ representative
for every $\alpha<1$, which agrees with the quasi-psh representative.

On a coordinate ball choose smooth $\rho$ with $\ddc\rho\geq\omega$.
The radial regularizations of the bounded psh function $u+\rho$ decrease
to it, and their complex Hessians converge uniformly on smaller balls.
Continuity of bounded products and the finite smooth-background
expansion identify $(\omega+\ddc u)^n$ with $\sigma^n$. For smooth real
$h$, the same argument applies to $u+th$, which is admissible for all
sufficiently small signed $t$ by uniform positivity. Define
\[
 F_t=F+\log\frac{(\sigma+t\ddc h)^n}{\sigma^n}.
\]
Its equation has the same normalization scalar as the equation for $u$,
so $B_0(F_t)=B_0(F)$. The supporting inequality gives
$\int_X(F_t-F)\,\mathrm d\nu\leq0$ for each
$\nu\in\mathcal A_0(F)$. Since
$F_t-F=tL_\sigma h+O(t^2)$ uniformly, using both signs yields stationarity.

Choose $p=4n+1$ and $q=p/(p-1)=1+1/(4n)$. By
\cref{cor:gaud-adjoint-lq}, every supporting probability has an $L^q$
density, since $q<n/(n-1)$. Smooth density in $W^{2,p}$ extends its
stationarity to every $W^{2,p}$ test. By
\cref{lem:continuous-hermitian-fredholm}, all these densities lie in a
one-dimensional space. Their integrals equal one, so they coincide.
Nonemptiness is \cref{thm:weak-adjoint-probabilities}, and the derivative
formula follows from \cref{prop:normalization-directional}.
\end{proof}

\subsection{Products with a Hermitian factor}

An explicit product construction exhibits the criterion for densities
with an arbitrary continuous modulus. Let $(Y,\kappa)$ be a connected
compact Hermitian $k$-fold, $k\geq1$, with $\int_Y\kappa^k=1$.
Denote by $\nu_\kappa$ the smooth stationary probability of
$L_\kappa=\operatorname{tr}_\kappa\ddc$, given by the Gauduchon multiplier
when $k\geq2$; for $k=1$ it is $\kappa$. Write $w$ for the
Fubini--Study area of mass one on $\mathbb P^1$. All factor forms and
functions below are pulled back to $X=\mathbb P^1\times Y$.

\begin{proposition}\label{prop:continuous-product-adjoint}
Let $a>0$, let $b$ be smooth and positive on $\mathbb P^1$, and let $A$
be continuous and positive there, with $\int_{\mathbb P^1}Aw=a$.
For
\[
 \omega=aw+b\kappa,\qquad \mathrm dV=w\wedge\kappa^k,
 \qquad F=\log((k+1)Ab^k),
\]
one has
\[
 c_0(F)=1,\qquad
 \mathcal A_0(F)=\left\{\frac Aa\,w\otimes\nu_\kappa\right\}.
\]
\end{proposition}
\begin{proof}
The Green solution of $\ddc u=(A-a)w$ on $\mathbb P^1$ has continuous
complex Hessian and is bounded. Its pullback yields
$\sigma=Aw+b\kappa>0$, and
$\sigma^{k+1}=(k+1)Ab^kw\wedge\kappa^k$.
Thus scalar uniqueness gives $c_0(F)=1$, and
\cref{thm:continuous-hessian-adjoint-uniqueness} gives a unique support.
For every smooth $h$ depending on all variables, the block decomposition
\[
 L_\sigma h=A^{-1}L_wh+b^{-1}L_\kappa h
\]
shows that $(A/a)w\otimes\nu_\kappa$ is stationary: the first term
integrates to zero by Stokes on $\mathbb P^1$, and the second by
stationarity on $Y$ and Fubini. This candidate has a bounded density
with respect to $\mathrm dV$. The one-dimensional annihilator of
\cref{lem:continuous-hermitian-fredholm} and probability normalization
therefore identify it with the supporting probability.
\end{proof}

\begin{example}\label{ex:non-dini-hopf-product-adjoint}
There is a non-K\"ahler threefold with a nonclosed Hermitian background
and a positive continuous density whose logarithm is neither Dini
continuous nor H\"older continuous of any positive order, but whose
supporting probability is explicit and unique.
\end{example}
\begin{proof}
Take the Hopf surface
$Y=(\mathbb C^2\setminus\{0\})/\langle z\mapsto2z\rangle$ and its
unit-volume metric $\kappa=c|z|^{-2}\ddc|z|^2$, where $c>0$ is chosen
so that $\int_Y\kappa^2=1$. In the notation of
\eqref{eq:hopf-torsion}, this is a multiple of $\alpha+\beta$.
Here $\alpha$ has rank one, so $\ddc\kappa=-c\alpha^2=0$ and
$\nu_\kappa=\kappa^2$. Polar coordinates identify $Y$ with
$S^3\times S^1$, hence $H^2(Y,\mathbb R)=0$. A closed positive form
would be exact and have positive square integral, contradicting Stokes.
Thus $Y$ is non-K\"ahler, and so is $\mathbb P^1\times Y$, by restriction.

Let $\tau\in[-1,1]$ be the height on $\mathbb P^1$, set $b=2+\tau$,
and put
\[
 f(t)=\begin{cases}
       1/\log(e/|t|),&0<|t|\leq1,\\
       0,&t=0,
      \end{cases}
 \qquad
 A=a+\varepsilon\left(f(\tau)-\int_{\mathbb P^1}f(\tau)w\right),
 \quad0<\varepsilon<a.
\]
Since $0\leq f\leq1$, this is a positive continuous function of mean
$a$. Along a normal segment to the equator, $|\tau|$ is comparable
to the distance $r$. The oscillation of $A$ is therefore bounded below
by a positive multiple of $1/\log(e/r)$. Its Dini integral
$\int_0^{r_0}\mathrm dr/(r\log(e/r))$ diverges, and this modulus is
larger than every $r^\gamma$, $\gamma>0$, as $r\downarrow0$.
The logarithm of $A$ has the same lower bound, because $A$ is uniformly
positive and bounded. The smooth term $2\log b$ contributes only $O(r)$,
so $F=\log(3Ab^2)$ has neither regularity. The mixed component
$\mathrm db\wedge\kappa$ shows that $\omega=aw+b\kappa$ is nonclosed.
By \cref{prop:continuous-product-adjoint}, its supporting probability
is $(A/a)w\wedge\kappa^2$.
\end{proof}

The theorem assumes a continuous positive complex Hessian for a bounded
solution. It does not establish that hypothesis for every continuous $F$;
the uniqueness question for arbitrary continuous data remains open.
